\documentclass[aspectratio=169]{beamer}
\usetheme{Madrid}
\usecolortheme{default}

\usepackage[utf8]{inputenc}
\usepackage[T1]{fontenc}
\usepackage{lmodern}
\usepackage{graphicx}
\usepackage{booktabs}
\usepackage{siunitx}

\title[Solara]{Solara: Smart Agricultural Monitoring\\and Crop Recommendation System}
\subtitle{Senior Design Presentation}
\author{Batuhan Y\i lmaz \and Muhammed Kaan Uman \and Emre Kemal Aksel}
\institute{Akdeniz University, Computer Engineering}
\date{May 2026}

\begin{document}

\begin{frame}
  \titlepage
\end{frame}

\begin{frame}{1) Problem: Why This Matters}
  \begin{itemize}
    \item Small and medium-scale farmers face an \textbf{information gap} in field decisions.
    \item Climate variability and resource costs make intuition-only farming risky.
    \item Over-irrigation and over-fertilization increase cost and damage soil/water health.
    \item Wrong crop selection for soil conditions leads to yield and income loss.
    \item Existing smart farming stacks are often too expensive for widespread adoption.
  \end{itemize}
\end{frame}

\begin{frame}{2) Objectives and Engineering Requirements}
  \textbf{Project objective:} Democratize precision agriculture with an affordable, data-driven platform.
  \vspace{0.5em}
  \begin{itemize}
    \item Build a low-cost ESP32 IoT node; target deep-field operation with cellular transport.
    \item Combine live telemetry with manual NPK/pH input (\textbf{hybrid data model}).
    \item Deliver top-N crop recommendations through a low-latency ML engine.
    \item Provide web, mobile, and admin interfaces for end-to-end field management.
  \end{itemize}
  \vspace{0.5em}
  \textbf{Constraints:}
  \begin{itemize}
    \item Cost-sensitive hardware choices
    \item Limited project timeline and deployment complexity
    \item Rural connectivity and battery-life limitations
  \end{itemize}
\end{frame}

\begin{frame}{3) Related Work: What Exists Today}
  \textbf{Typical patterns in recent studies}
  \begin{itemize}
    \item IoT + ML crop recommendation pipelines are common and effective for decision support.
    \item Many works show strong prediction quality on controlled datasets.
    \item Mobile/cloud dashboards are widely used for farmer-facing access.
    \item Some studies include explainability concepts, but usually at a conceptual level.
  \end{itemize}
  \vspace{0.5em}
  \textbf{Representative directions from your review}
  \begin{itemize}
    \item Real-time monitoring + recommendation apps
    \item Nutrient-focused sensing systems
    \item Time-series deep models for yield/crop decisions
    \item Cloud-first precision agriculture platforms
  \end{itemize}
\end{frame}

\begin{frame}{4) Related Work: Main Limitations}
  \textbf{Economic barrier}
  \begin{itemize}
    \item Several systems assume automated NPK/pH probes, which are costly for small farms.
  \end{itemize}
  \textbf{Infrastructure barrier}
  \begin{itemize}
    \item Approaches often depend on reliable local connectivity or extra gateway setup.
  \end{itemize}
  \textbf{Usability barrier}
  \begin{itemize}
    \item Outputs are often raw predictions; guidance/explanation for farmers is limited.
  \end{itemize}
  \textbf{Integration barrier}
  \begin{itemize}
    \item Many solutions optimize one layer (sensor, model, or app) rather than full workflow.
  \end{itemize}
\end{frame}

\begin{frame}{5) Solara Positioning and Novelty}
  \textbf{How Solara addresses those gaps}
  \begin{itemize}
    \item \textbf{Hybrid affordability:} manual NPK/pH + automated telemetry to reduce hardware cost.
    \item \textbf{End-to-end architecture:} firmware, broker, backend, ML service, web/mobile/admin.
    \item \textbf{Operational analysis modes:} range, future baseline, and what-if scenario support.
    \item \textbf{Guided interaction:} crop-guide-driven chatbot for practical Q\&A.
    \item \textbf{Code reality note:} current repository firmware publishes via Wi-Fi MQTT; transport is abstracted for deep-field extensions.
  \end{itemize}
\end{frame}

\begin{frame}{6) Solution Overview (System Architecture)}
  \begin{itemize}
    \item \textbf{Perception layer:} ESP32 + sensors (BME280, DS18B20, capacitive moisture)
    \item \textbf{Transport layer:} MQTT via Mosquitto (current firmware publishes over Wi-Fi)
    \item \textbf{Processing layer:} Spring Boot + FastAPI ML engine + PostgreSQL/PostGIS
    \item \textbf{Interaction layer:} React web panel, React Native mobile app, admin panel
  \end{itemize}
  \vspace{0.5em}
  \centering
  \fbox{\parbox{0.9\linewidth}{\centering Insert high-level architecture diagram here\\(Report Figure 3.4)}}
\end{frame}

\begin{frame}{7) Key Design Decisions and Trade-offs}
  \begin{itemize}
    \item \textbf{LightGBM over alternatives:} better tabular efficiency, low-latency inference.
    \item \textbf{Transport trade-off:} current code uses Wi-Fi MQTT for stable development; GSM path is a deep-field extension with higher energy cost.
    \item \textbf{Hybrid data model:} lower cost; requires periodic manual NPK/pH entry.
    \item \textbf{Cloud inference:} simpler updates/scaling; requires network availability.
    \item \textbf{Deep sleep firmware:} extends battery life with scheduled transmission windows.
  \end{itemize}
\end{frame}

\begin{frame}{8) Implementation Snapshot}
  \begin{itemize}
    \item Secure auth and email verification with JWT + SMTP.
    \item Field mapping and geospatial storage using PostGIS polygons.
    \item Device pairing via MAC snippet and whitelist-based registry.
    \item Historical charts and telemetry cards for 7/14/30 day insights.
    \item Three recommendation scenarios:
    \begin{itemize}
      \item Range analysis (historical observation)
      \item Future season planning (5-year baseline)
      \item What-if simulation (manual environmental overrides)
    \end{itemize}
  \end{itemize}
\end{frame}

\begin{frame}{9) Demo Plan (Live if Possible)}
  \textbf{Target flow (2--3 minutes)}
  \begin{enumerate}
    \item Create/select a field and show map polygon
    \item Pair device (or load already paired field)
    \item Show live telemetry and trend chart
    \item Run crop recommendation and show ranked outputs
    \item Ask chatbot a crop-guide question
  \end{enumerate}
  \vspace{0.5em}
  \textbf{Backup plan:} pre-recorded video + screenshots from web/mobile.
\end{frame}

\begin{frame}{10) Results: ML and System Performance}
  \begin{itemize}
    \item \textbf{LightGBM accuracy:} \textbf{99.39\%} on validation and test splits.
    \item Ranked top-N output improves decision support over single-label output.
    \item Feature analysis highlights rainfall, humidity, and potassium importance.
    \item Reliable telemetry ingestion over MQTT in current stack (backend subscriber + persisted logs).
  \end{itemize}
  \vspace{0.5em}
  \centering
  \fbox{\parbox{0.9\linewidth}{\centering Insert confusion matrix/feature importance\\(Report Figure 4.1)}}
\end{frame}

\begin{frame}{11) Results: Reliability and Scalability}
  \begin{itemize}
    \item Sensor telemetry pipeline stores ambient/soil/battery/location payload fields end-to-end.
    \item Deep sleep workflow is implemented in firmware for periodic reporting cycles.
    \item PostGIS-backed field polygons are enabled through geometry(Polygon, 4326) modeling.
    \item CSV export for telemetry and analysis log persistence are implemented in backend APIs.
    \item Unauthorized or unpaired device IDs are ignored during ingestion.
  \end{itemize}
\end{frame}

\begin{frame}{12) Challenges and Lessons Learned}
  \begin{itemize}
    \item Balancing affordability with data completeness required a hybrid input strategy.
    \item Rural connectivity constraints drove protocol and power-management decisions.
    \item Cross-platform synchronization (web, mobile, admin, ML, chatbot) needed modular APIs.
    \item Explainability matters: users need guidance, not only model scores.
    \item Dataset origin limits localization; model quality still depends on regional data.
  \end{itemize}
\end{frame}

\begin{frame}{13) Conclusion and Impact}
  \begin{itemize}
    \item Solara delivers an end-to-end precision agriculture platform for cost-sensitive users.
    \item The project validates a practical path from raw field data to actionable recommendations.
    \item Hybrid data collection lowers barrier to entry without removing intelligence.
    \item Architecture is production-oriented: modular, containerized, and scalable.
  \end{itemize}
\end{frame}

\begin{frame}{14) Future Work}
  \begin{itemize}
    \item Train with localized Turkish/Mediterranean agronomic datasets.
    \item Add offline-first mobile support with local cache and sync.
    \item Explore TinyML anomaly detection on-device for smarter transmission control.
    \item Integrate market/economic APIs for profit-aware recommendations.
    \item Extend from advisory to closed-loop irrigation/fertigation automation.
  \end{itemize}
\end{frame}

\begin{frame}{Q\&A}
  \centering
  \Large Thank you.\\
  \vspace{1em}
  Questions?
\end{frame}

\begin{frame}{Backup: Why LightGBM?}
  \begin{itemize}
    \item Efficient on tabular agricultural features.
    \item Leaf-wise growth and GOSS improve speed/accuracy trade-off.
    \item Supports multiclass prediction with ranked probability outputs.
    \item Suitable for low-latency cloud inference in mobile/web workflows.
  \end{itemize}
\end{frame}

\begin{frame}{Backup: Main Limitations}
  \begin{itemize}
    \item Current model trained on Kaggle data, not yet locally specialized.
    \item Manual NPK/pH entry introduces user dependency and possible input errors.
    \item Firmware in repository currently uses Wi-Fi MQTT; GSM production hardening remains pending.
    \item Hardware prototype is monitoring-first, not yet automatic actuation.
  \end{itemize}
\end{frame}

\end{document}
